% Pursuing Stacks (cotes 134-2 to 134-8): archive cote 134-3, site page 63 (PDF page 64), typescript page 63.
% typescript. Draft compared: ps2.tex:3244-3313.
% AI-transcribed by Claude Opus 5.5 (claude-opus-5-5, adaptive thinking, effort high), Message Batches API, batch ps1 (msgbatch_01Rv5hCaNTtUt4JHx1yqoHPL), 2026-11-12; not yet checked by eye.
\scanpage{134-3/63}
\mspage{63}

22.3.

38. Lets \inserted{get} back to the ``asphericity story retold'' - I had to stop yester-%
day just in the middle, as it was getting prohibitively late.

I want to comment a little about the ``key result'' just stated and
proved. T\hand{h}e main point of this result, forgetting the game of giving
heaps of equivalent formulations of two kind of properties, is that the
extremely simple condition b (iii), namely that $i^{*}(\hand{$\Delta$}_1)$ is aspheric over
the unit element $e_{A^\wedge}$ of A,\inserted{plus asphericity of the latter,} ensure\deleted{s} already the conditions in a), which
can be viewed (among others) as just stating that the compositum
\[ i_A i^{*} : \text{(Cat)} \hand{$\longrightarrow$} A^\wedge \hand{$\longrightarrow$} \text{(Cat)} \]
\xout{induces} from (Cat) to (Cat) induces an autoequivalence of the localised ca%
tegory (Hot), or (what amounts still to the same) a functor (Hot)\hand{$\to$}(Hot)
\xout{equi} isomorphic to the identity functor. (NB that the two do indeed amount
to the same follows at once from the implication (iv)\hand{$\Rightarrow$}(ii) in a).) It
is interesting to note that both properties, the stronger one that
$i^{*}(\hand{$\Delta$}_1)$ is aspheric over $e_{A^\wedge}$, and the weaker one in terms of properties
of the compositum $i_A i^{*}$ , \uline{make a sense without any reference to the extra}
\uline{assumption that the categories} i(a) \uline{have a final element each}, nor to the
corresponding map ($*$) $A_{/C}$ \hand{$\to$} C (functorial in C). This suggests that
there should exist a more general statement than in the theorem, without
making the assumption about final objects in the categories \over{a}{i}(a), and
without the possibility of a \uline{direct} comparison of $A_{/C}$ and C through a
functor between them. Indeed I have an idea of a statement in this res-%
\margin{X} pect, however for the time being it seems that the theorem as stated is
sufficiently general for handling the situations I have in mind \hand{(44)}.
\margin{cf 65 E,F (after ``\unsure{prousin}'' 47)}

Applying the theorem to the \uline{canonical} functor $i$
\[ i_o : a \hand{$\mapsto$} A_{/a} \ : \ A \hand{$\longrightarrow$} \text{(Cat)} , \]
whose canonical extension to $A^\wedge$ (as a functor $A^\wedge$\hand{$\to$}(Cat) commuting to
direct limits) is the functor \deleted{\illegible{1}}
\[ i_A : F \hand{$\to$} A_{/F} \ \hand{\unsure{;|:}} \ A^\wedge \hand{$\longrightarrow$} \text{(Cat)} , \]
giving rise to the right adjoint
\[ i_o^{*} = j_A : \over{A}{(}\text{Cat)} \hand{$\longrightarrow$} A^\wedge , \]
the condition (i) in a) is nothing but the familiar condition
\[ i_A j_A(C) \hand{$\to$} C \quad \text{\inserted{a weak equivalence} \deleted{aspheric} for any C in (Cat)} , \]
which we had used for our first (or rather, second already !) définition
of so-called``test categories''. Later on we considerably strengthened this
condition - we now call them ``test categories in the wide sense''. On \xout{th}
the other hand, as we already noticed before, here $i_o^{*}(\hand{$\Delta$}_1) = j_A(\hand{$\Delta$}_1)$ \uline{is}
\uline{nothing but the Lawvere element} $L_A$ \uline{of} $A^\wedge$. Thus the main content of the
